\documentclass[11pt,a4paper]{article}
\usepackage[utf8]{inputenc}
\usepackage[french]{babel}
\usepackage[T1]{fontenc}
\usepackage{lmodern}
\usepackage[margin=1.5cm]{geometry}

\begin{document}

% En-tête
\begin{center}
    {\Huge \textbf{Hugo Leroy}} \\[0.2cm]
    {\Large \textbf{Ingénieur IA}} \\[0.1cm]
    hugo.leroy@email.com \ | \ 06 13 89 54 58
\end{center}

\vspace{0.3cm}

% Résumé / Profil
\section*{Profil Professionnel}
\noindent
Expert en Data Science et Intelligence Artificielle, je conçois des modèles prédictifs performants en exploitant des volumes massifs de données (Big Data). Maîtrisant les algorithmes de Machine Learning et de Deep Learning, j'accompagne les entreprises dans l'automatisation de leurs processus métiers et la valorisation de leur patrimoine de données, tout en assurant une mise en production fiable des algorithmes développés.

\vspace{0.3cm}

% Compétences
\section*{Compétences Techniques}
\noindent
\textbf{PyTorch} $\bullet$ \textbf{Pandas} $\bullet$ \textbf{Hadoop} $\bullet$ \textbf{SQL} $\bullet$ \textbf{Machine Learning} $\bullet$ \textbf{Scikit-Learn} $\bullet$ \textbf{Deep Learning} $\bullet$ \textbf{Databricks}

\vspace{0.3cm}

% Expériences
\section*{Expériences Professionnelles}
\noindent \textbf{Ingénieur IA / Confirmé} --- \textit{Dataiku} \hfill \textbf{10/2024 à Aujourd'hui} \\
Optimisation des hyperparamètres de modèles Random Forest réduisant le taux de faux positifs de 22\% dans le système de détection de fraudes. Conception d'un modèle de prévision de la demande (Time Series) avec Prophet, réduisant les ruptures de stock de 30\% sur l'année. Conception d'un système de recommandation basé sur le Deep Learning, générant une hausse de 15\% du chiffre d'affaires sur le segment e-commerce. Collaboration avec les experts métiers pour définir les KPIs de succès et extraire les features les plus pertinentes des bases de données historiques. \\[0.3cm]
\noindent \textbf{Ingénieur IA} --- \textit{Dataiku} \hfill \textbf{10/2023 à 09/2024} \\
Déploiement de modèles prédictifs en production via FastAPI et Docker, avec intégration complète d'une boucle de feedback utilisateur. Conception d'une architecture MLOps robuste sur AWS SageMaker pour garantir le suivi des performances et le réentraînement automatique des algorithmes. Mise en place d'un pipeline Big Data avec Apache Spark pour traiter et nettoyer plus de 500 To de données par jour en temps réel. Entraînement et déploiement de modèles de Machine Learning (NLP) ayant permis d'automatiser 60\% des tâches manuelles de classification de documents. \\[0.3cm]
\noindent \textbf{Stage — Assistant Ingénieur IA} --- \textit{Artefact} \hfill \textbf{07/2022 à 10/2023} \\
Mise en place d'une architecture de Feature Store centralisée, facilitant la réutilisation des variables entre les différentes équipes Data. Création d'un algorithme de clustering non supervisé pour segmenter la base client en 5 catégories distinctes et optimiser les campagnes marketing. Analyse exploratoire des données (EDA) et création de dashboards interactifs avec Tableau, permettant au comité de direction de piloter l'activité. \\[0.3cm]


\vspace{0.3cm}

% Formations
\section*{Formation \& Diplômes}
\noindent \textbf{Master en Data Science} \hfill \textbf{2021 - 2023} \\
\textit{Claude Bernard Lyon 1} \\[0.3cm]
\noindent \textbf{Licence en Data Science} \hfill \textbf{2018 - 2021} \\
\textit{Claude Bernard Lyon 1} \\[0.3cm]


\end{document}